\documentclass[11pt]{article}
\usepackage[a4paper,margin=0.65in]{geometry}
\usepackage{booktabs,longtable,array,pdflscape}
\usepackage{url,hyperref,microtype}
\newcolumntype{L}[1]{>{\raggedright\arraybackslash\hspace{0pt}}p{#1}}
\setlength\LTleft{0pt}
\setlength\LTright{0pt}
\emergencystretch=3em
\sloppy

\title{Supplementary Table 10. TRIPOD+AI reporting checklist mapping for the manuscript}
\author{Frontiers in Signal Processing R1 supplementary table}
\date{07 July 2026}

\begin{document}
\maketitle

\section*{Supplementary Table 10. TRIPOD+AI reporting checklist mapping for the manuscript}

Standalone display version for LaTeX upload/review. Machine-readable source: source\_\allowbreak{}tables/\allowbreak{}Supplementary\_\allowbreak{}Table\_\allowbreak{}10\_\allowbreak{}TRIPOD\_\allowbreak{}AI\_\allowbreak{}checklist.\allowbreak{}csv. Rows: 52; columns: 5. The display below contains the complete source table.

\begin{landscape}
\scriptsize
\setlength\LTleft{0.025\linewidth}
\setlength\LTright{0.025\linewidth}
\setlength{\tabcolsep}{2pt}
\renewcommand{\arraystretch}{1.16}
\begin{longtable}{L{0.184\linewidth} L{0.184\linewidth} L{0.184\linewidth} L{0.184\linewidth} L{0.184\linewidth}}
\caption{Supplementary Table 10}\label{tab:standalone_s10_g1_block1}\\
\toprule
item & topic & status & reported location & comment \\
\midrule
\endfirsthead
\multicolumn{5}{l}{\textit{Continued from Supplementary Table 10.}}\\
\toprule
item & topic & status & reported location & comment \\
\midrule
\endhead
\endfoot
\bottomrule
\endlastfoot
1 & Title & Reported & Title & Identifies self-\allowbreak{}supervised CTG trajectory representation and calibrated neonatal risk enrichment in the target intrapartum CTG population.\allowbreak{} \\
2 & Abstract & Reported & Structured abstract & Structured Background Objectives Methods Results Conclusions abstract aligned with Frontiers in Signal Processing positioning.\allowbreak{} \\
3a & Healthcare context and rationale & Reported & Introduction & Explains intrapartum CTG decision-\allowbreak{}support context and limitations of static or subjective interpretation.\allowbreak{} \\
3b & Target population and intended purpose & Reported & Introduction and Discussion & Target population is intrapartum CTG records;\allowbreak{} intended use is risk enrichment and phenotype interpretation rather than autonomous diagnosis.\allowbreak{} \\
3c & Known health inequalities & Partially reported & Discussion limitations & No protected sociodemographic fairness analysis was possible beyond available CTU-\allowbreak{}UHB metadata;\allowbreak{} this is treated as a limitation.\allowbreak{} \\
4 & Objectives & Reported & Introduction & Objective is evaluation of an open-\allowbreak{}data self-\allowbreak{}supervised CTG trajectory-\allowbreak{}representation workflow for calibrated neonatal risk enrichment.\allowbreak{} \\
5a & Data sources & Reported & Methods and Supplementary Table 1 & CTU-\allowbreak{}UHB/\allowbreak{}CTU-\allowbreak{}CHB is primary raw-\allowbreak{}signal cohort;\allowbreak{} CTGDL-\allowbreak{}FHRMA is external representation/\allowbreak{}domain-\allowbreak{}shift resource;\allowbreak{} UCI CTG is feature-\allowbreak{}level reference.\allowbreak{} \\
5b & Data dates & Partially reported & Supplementary Table 1 & Public datasets are cited with access information;\allowbreak{} detailed original accrual dates depend on source metadata.\allowbreak{} \\
6a & Study setting & Reported & Methods and Supplementary Table 1 & Intrapartum cardiotocography public data from hospital-\allowbreak{}based CTG resources.\allowbreak{} \\
6b & Eligibility criteria & Reported & Methods and Supplementary Table 1 & Records included when raw FHR/\allowbreak{}UC and required metadata were available for the analysis pipeline.\allowbreak{} \\
6c & Treatments received & Not applicable & Methods & No treatment-\allowbreak{}effect modelling was performed;\allowbreak{} delivery process variables were not used as intervention exposures.\allowbreak{} \\
7 & Data preparation and quality checking & Reported & Methods & Signal windowing, missingness, feature extraction, train-\allowbreak{}only scaling and source-\allowbreak{}data generation are described.\allowbreak{} \\
8a & Outcome definition and time horizon & Reported & Methods outcome definition & Primary endpoint is cord pH <7.\allowbreak{}15 or 5-\allowbreak{}minute Apgar <7 at record level.\allowbreak{} \\
8b & Subjective outcome assessment & Not applicable & Methods outcome definition & Outcome uses recorded pH and Apgar metadata from the public dataset rather than new subjective adjudication.\allowbreak{} \\
8c & Outcome blinding & Not applicable & Methods & Retrospective public-\allowbreak{}data analysis;\allowbreak{} outcome labels were withheld from SSL pretraining.\allowbreak{} \\
9a & Initial predictor choice & Reported & Methods outcome and prediction models & Classical signal features and SSL embeddings were prespecified from CTG physiology and self-\allowbreak{}supervised representation design.\allowbreak{} \\
9b & Predictor definitions & Reported & Methods and supplementary tables & Signal features, SSL embedding dimensions and phenotype variables are described;\allowbreak{} source-\allowbreak{}data tables are included.\allowbreak{} \\
9c & Subjective predictor assessment & Not applicable & Methods & Predictors were algorithmically extracted from public CTG signals and metadata.\allowbreak{} \\
10 & Sample size & Reported & Methods and Results & Record counts, event counts and window counts are reported;\allowbreak{} limitations discuss small event count and high-\allowbreak{}dimensional model risk.\allowbreak{} \\
11 & Missing data & Reported & Methods and Results & Median imputation, missingness features and no-\allowbreak{}missingness sensitivity are reported.\allowbreak{} \\
12a & Data partitioning and analysis use & Reported & Methods record split & Record-\allowbreak{}level training/\allowbreak{}test split, train-\allowbreak{}only SSL pretraining and held-\allowbreak{}out validation are stated.\allowbreak{} \\
12b & Predictor handling & Reported & Methods & Numeric predictors were median-\allowbreak{}imputed and standardized;\allowbreak{} categorical phenotypes were one-\allowbreak{}hot encoded.\allowbreak{} \\
12c & Model type and tuning & Reported & Methods and Supplementary Table 9 & Transformer SSL, logistic pipelines, exploratory compact sweep and development-\allowbreak{}only selection sensitivity are described.\allowbreak{} \\
12d & Heterogeneity across clusters & Partially reported & Results and Discussion & Domain-\allowbreak{}shift analysis compares CTU-\allowbreak{}UHB with CTGDL-\allowbreak{}FHRMA;\allowbreak{} formal multi-\allowbreak{}centre heterogeneity was not possible.\allowbreak{} \\
12e & Performance measures & Reported & Methods and Results & AUROC, AUPRC, Brier score, calibration, decision curve and bootstrap CIs are reported with rationale.\allowbreak{} \\
12f & Model updating or recalibration & Reported & Methods and Results & Cross-\allowbreak{}fitted Platt and isotonic recalibration are described;\allowbreak{} Platt scaling is emphasized.\allowbreak{} \\
12g & Prediction calculation & Reported & Methods and Code availability & Pipeline scripts are named and the versioned repository archive is available at https:\allowbreak{}/\allowbreak{}/\allowbreak{}doi.\allowbreak{}org/\allowbreak{}10.\allowbreak{}5281/\allowbreak{}zenodo.\allowbreak{}21242760.\allowbreak{} \\
13 & Class imbalance & Reported & Methods and Results & Class-\allowbreak{}balanced logistic regression and AUPRC emphasis are justified under event imbalance.\allowbreak{} \\
14 & Fairness & Partially reported & Discussion limitations & Available public metadata did not support a full fairness assessment;\allowbreak{} age and clinical stratification were limited.\allowbreak{} \\
15 & Model output & Reported & Methods and Discussion & Outputs are risk-\allowbreak{}enrichment scores and calibrated probabilities;\allowbreak{} not intended as definitive diagnosis.\allowbreak{} \\
16 & Training versus evaluation differences & Reported & Methods and Results & CTU train/\allowbreak{}test split and CTGDL representation-\allowbreak{}domain differences are explicitly described.\allowbreak{} \\
17 & Ethical approval & Reported & Declarations & Public de-\allowbreak{}identified datasets;\allowbreak{} ethics approval not applicable.\allowbreak{} \\
18a & Funding & Reported & Declarations & No specific funding.\allowbreak{} \\
18b & Conflicts of interest & Reported & Declarations & No competing interests declared.\allowbreak{} \\
18c & Protocol & Reported & Code availability and limitations & No formal protocol was registered;\allowbreak{} analysis scripts and source data are included for reproducibility.\allowbreak{} \\
18d & Registration & Reported & Code availability and limitations & Study was not registered.\allowbreak{} \\
18e & Data sharing & Reported & Availability of data and materials & Public data sources and derived source-\allowbreak{}data tables are listed.\allowbreak{} \\
18f & Code sharing & Reported & Code availability & Local scripts are identified;\allowbreak{} the GitHub repository and versioned Zenodo archive are reported in Code availability.\allowbreak{} \\
19 & Patient and public involvement & Not applicable & Methods/\allowbreak{}Declarations & No patient or public involvement in this retrospective public-\allowbreak{}data analysis.\allowbreak{} \\
20a & Participant flow & Reported & Results and Figure 1 & Cohort scale and train/\allowbreak{}test records are shown.\allowbreak{} \\
20b & Participant characteristics & Reported & Supplementary Table 1 & Dataset scale and available record-\allowbreak{}level characteristics are summarized in supplementary files.\allowbreak{} \\
20c & Development versus evaluation comparison & Reported & Figure 1 and Supplementary Table 1 & Train/\allowbreak{}test event balance and external domain-\allowbreak{}shift analysis are reported.\allowbreak{} \\
21 & Number of participants and events per analysis & Reported & Results and Supplementary Table 9 & Events are reported for primary and sensitivity analyses.\allowbreak{} \\
22 & Model specification & Reported & Methods and Code availability & Architecture, feature dimensions and logistic models are described;\allowbreak{} implementation scripts are included in the GitHub repository and versioned Zenodo archive.\allowbreak{} \\
23a & Model performance with uncertainty & Reported & Results and Supplementary Tables 2-\allowbreak{}3 & Performance estimates include bootstrap confidence intervals and model differences.\allowbreak{} \\
23b & Heterogeneity in performance & Partially reported & Results and Discussion & Domain-\allowbreak{}shift and missingness sensitivity are reported;\allowbreak{} subgroup fairness analysis remains limited.\allowbreak{} \\
24 & Model updating results & Reported & Results and Supplementary Table 7 & Platt and isotonic recalibration results are reported.\allowbreak{} \\
25 & Interpretation & Reported & Discussion & Interprets risk enrichment, phenotype physiology, calibration and external-\allowbreak{}domain boundaries.\allowbreak{} \\
26 & Limitations & Reported & Discussion & Discusses retrospective design, model-\allowbreak{}selection sensitivity, event count, dimensionality, missingness and lack of external outcome validation.\allowbreak{} \\
27a & Handling poor quality or unavailable inputs & Partially reported & Methods and Discussion & Missingness sensitivity and signal-\allowbreak{}quality features are reported;\allowbreak{} deployment handling is future work.\allowbreak{} \\
27b & User interaction and expertise & Reported & Discussion & Model is positioned for clinician-\allowbreak{}facing decision support rather than autonomous use.\allowbreak{} \\
27c & Future research & Reported & Discussion & Calls for larger locked train/\allowbreak{}development/\allowbreak{}test designs, external outcome validation and prospective evaluation.\allowbreak{} \\
\end{longtable}
\normalsize
\end{landscape}


\end{document}
